% BEGIN R28_MEMORY_PARITY
% Additive continuation of MEM01: both readers are already defined.
\begin{proposition}[Parity of the readers under oriented reversal]
With cyclic indices in \(\mathbb Z/12\mathbb Z\), let
\((C_M\tau)_m=-\tau_{-m}\) and \(\operatorname{sgn}(0)=0\).
The two readers satisfy
\[
 \nu_{120}(C_M\tau)=-\nu_{120}(\tau),\qquad
 \nu_{270}(C_M\tau)=\nu_{270}(\tau).
\]
\end{proposition}
\begin{proof}
The reflection \(m\mapsto-m\) permutes the four position classes
modulo four. Reversal changes the sign of each class sum and hence,
by the oddness of \(\operatorname{sgn}\), the sign of \(\nu_{120}\).
In the second reader, the signs of the two factors cancel:
\[
 \sum_m(C_M\tau)_m(C_M\tau)_{m+3}
 =\sum_m\tau_{-m}\tau_{-m-3}
 =\sum_k\tau_{k+3}\tau_k
 =\nu_{270}(\tau).
\]
\end{proof}
Orientation thus distinguishes an odd reading from an even correlation.
Both depend on the arrangement of samples in the record; their
transformation laws specify which information persists when the route
is reversed.
% END R28_MEMORY_PARITY
